\scanpage{289}
\mspage{239}
\diagram{\begin{tikzcd}[ampersand replacement=\&]
 \unsure{\mathscr X|\mathcal X}_{U\cap U'} \arrow[r] \arrow[d,"{\varphi_1}"] \& \mathscr X_U \arrow[d,"{\psi_1}"]\\
 \mathscr X_{U'} \arrow[r] \& \mathscr X_V
\end{tikzcd}}
Donc par $W$(7 bis), (cas dual \deleted{\unsure{d}} des imm. fermées), on a $\varphi_1\in W\Longleftrightarrow\psi_1\in W$, qed.

Notons pour mémoire le corollaire qui nous avait motivé

\uline{Corollaire} Supposons $W$(7 bis) \inserted{(comme dans \unsure{tte} cette section)}. Soit un carré cocartésien dans $Cat$
\diagram{\begin{tikzcd}[ampersand replacement=\&]
 X_0 \arrow[r,"{i_1}"] \arrow[d,"{i_2}"'] \& X_1 \arrow[d]\\
 X_2 \arrow[r] \& X_3,
\end{tikzcd}}
avec $i_1$ (ou $i_2$) une $W$-cofibration (p. ex. une immersion de Dwyer). Alors ce carré est $W$-cocartésien.

Cette digression (p. 232 ff) n'est pas \unsure{nécessaire} d'ailleurs pour formuler dès maintenant la conclusion de \unsure{nos|mes} cogitations sur les cofibres $W$-homotopiques : \uline{Il y a une belle théorie dans} $Cat$, pourvu seulement que $W$ satisfasse à $W$(7 bis). En effet \unsure{tout} résulte du fait que \unsure{toute} flèche \inserted{$f:X\longrightarrow Y$} dans $Cat$ se factorise en
\[
 X\xrightarrow{i_f\deleted{\unsure{p}}}\widetilde Y\xrightarrow{\unsure{p_f|\rho_f}}Y
\]
avec $i_f$ une $W$-cofibration (et $=$ une imm

% archive sheet number: 289

